% =====================================================================
%  Volatility Forecasting — Replication of CSV (2023, JFE)
%  Three-page (body) coursework report, two-column 9pt for density.
%  Tables and figures live in the appendix — they do not count toward
%  the three-page body limit.
% =====================================================================
\documentclass[twocolumn,9pt]{extarticle}

% ---------- geometry / typography ----------
\usepackage[a4paper,margin=1.5cm,columnsep=0.55cm]{geometry}
\usepackage{setspace}
\setstretch{0.92}
\usepackage[T1]{fontenc}
\usepackage[utf8]{inputenc}
\usepackage{lmodern}
\usepackage{microtype}
\usepackage{amsmath,amssymb}
\usepackage{textcomp}
\usepackage{booktabs}
\usepackage{graphicx}
\usepackage{array}
\usepackage{multirow}
\usepackage{caption}
\usepackage{float}
\usepackage[hidelinks]{hyperref}

% Tight section / paragraph spacing
\usepackage{titlesec}
\titlespacing*{\section}{0pt}{4pt}{2pt}
\titlespacing*{\subsection}{0pt}{3pt}{1pt}
\titleformat{\section}{\normalsize\bfseries}{\thesection.}{0.5em}{}
\titleformat{\subsection}{\small\bfseries}{\thesubsection}{0.5em}{}
\captionsetup{font=footnotesize,labelfont=bf,skip=2pt}
\setlength{\parskip}{1.5pt}
\setlength{\parindent}{0pt}
\setlength{\abovedisplayskip}{2pt}
\setlength{\belowdisplayskip}{2pt}

% Lead-in style "Question." "What CSV do." bold inline
\newcommand{\lead}[1]{\textbf{#1}\hspace{0.25em}\ignorespaces}

% ---------- title block ----------
\title{\vspace{-1.8em}\bfseries A Machine Learning Approach to Volatility Forecasting:\\
\normalsize Replication of Christensen, Siggaard \& Veliyev (2023, JFE)}
\author{\small Financial Analytics and Machine Learning — Coursework Report}
\date{\vspace{-1em}}

\begin{document}

% =====================================================================
%   COVER PAGE  (place file ucl_logo.png in project root)
% =====================================================================
\begin{titlepage}
\centering
\vspace*{2cm}

\includegraphics[width=6cm,trim={6pt 6pt 6pt 6pt},clip]{ucl_logo.png}\\[2cm]

\rule{\linewidth}{0.6pt}\\[0.4em]
{\Large\scshape University College London}\\[0.4em]
{\normalsize\scshape Institute of Finance and Technology}\\[0.3em]
\rule{\linewidth}{0.6pt}\\[3cm]

{\Huge\bfseries A Machine Learning Approach to\\[0.35em]
Volatility Forecasting\par}
\vspace{1.2em}
{\LARGE Replication of Christensen, Siggaard \& Veliyev\\(2023, Journal of Financial Econometrics)\par}
\vspace{3cm}

\begin{tabular}{rl}
\textbf{Student ID:}      & 25130184 \\[0.35em]
\textbf{Name:}            & Shakhzod Makhmudov \\[0.35em]
\textbf{Module:}          & IFTE0004 \\
\end{tabular}

\vfill
{\normalsize Code repository: \url{https://github.com/Shakhzod555/Financial-Analytics-and-Machine-Learning-.git}}
\vspace{1cm}
\end{titlepage}

% Restart page numbering at 1 for the main body
\setcounter{page}{1}

% =====================================================================
%   MAIN BODY (twocolumn resumes automatically after titlepage)
% =====================================================================
\thispagestyle{empty}

% =====================================================================
\section{The paper}
% =====================================================================

\lead{Question.} Volatility forecasts are central to risk management
(VaR under Basel III's standardised market-risk charge), option
pricing (variance-swap valuation, gamma hedging) and portfolio
allocation. Realised variance (RV), the sum of squared intraday
returns, is the standard latent-volatility proxy (Andersen \&
Bollerslev, 1998). RV's defining feature is long memory:
autocorrelations decay hyperbolically rather than exponentially.
Corsi's (2009) Heterogeneous Autoregressive (HAR) model captures
that decay parsimoniously with three lagged-RV averages (daily,
weekly, monthly) fit by OLS. HAR survives as the benchmark because
its parsimony protects against overfit. That parsimony is the
puzzle. Christensen, Siggaard \& Veliyev (2023, henceforth CSV) ask
two linked questions: can machine learning extract more predictive
content than HAR's three-lag approximation, and what mechanism
delivers any gain.

\lead{What CSV do.} CSV position their paper as the first systematic
horse-race of machine learning against the HAR-class benchmark for
realised-variance forecasting at multiple horizons, with mechanism
analysis via Accumulated Local Effect plots (Apley \& Zhu, 2020) and
a Value-at-Risk application as economic validation. The horse-race
runs 22 models on 29 DJIA constituents, 2001--2017
($T = 4{,}257$ daily observations per stock). The HAR lineage
(six variants: HAR, HAR-X, LogHAR, LevHAR, SHAR, HARQ, the last two
from Patton \& Sheppard 2015 and Bollerslev, Patton \& Quaedvlieg
2016) is set against five regularised linears (Ridge, Lasso, Elastic
Net, Adaptive Lasso, Post-Lasso; Tibshirani 1996, Zou \& Hastie 2005),
three tree ensembles (Bagging, Random Forest, Gradient Boosting;
Breiman 1996, 2001, Friedman 2001), and four feed-forward neural
networks (NN$_1$--NN$_4$ in single-fit and 10-of-100 ensemble form).
Two information sets are tested: M\_HAR (the three HAR lags only) and
M\_ALL (HAR lags plus nine covariates including implied volatility,
VIX, ADS, EPU [Baker, Bloom \& Davis, 2016] and dollar volume).
Forecast comparison uses Diebold--Mariano (1995) with the
Harvey--Leybourne--Newbold (1997) small-sample correction and the
Model Confidence Set of Hansen, Lunde \& Nason (2011). CSV present
five findings. (i) On M\_HAR, deeper neural networks beat HAR by
roughly 5\,\% on average, regularised linears match HAR, Gradient
Boosting is slightly worse and Bagging materially worse
($\sim15\,\%$ above HAR at $h=1$). (ii) On M\_ALL, regularisation
cuts MSE by 6--9\,\% (Elastic Net at 8.4\,\%) and Random Forest by
9.9\,\%, while unregularised extended HAR overfits. (iii) ALE
attributes most predictive content to RVD, RVW and implied
volatility, with rankings differing across models. (iv) Gains rise
with horizon: at $h=22$ the Random Forest reduces MSE by roughly
40\,\% against HAR on M\_ALL, an effect CSV attribute to ML
approximating RV's long-memory persistence, with the slower ACF
decay of fitted values (their Figure~8) as supporting evidence.
(v) A closing Value-at-Risk application yields the same qualitative
ranking with smaller and less statistically significant gains.

\lead{Open questions.} Four weaknesses in CSV motivate the analysis
here. First, the long-horizon mechanism is identified only via the
ACF of fitted RV, a pattern consistent with both genuine long memory
and infrequent regime-switching with similar autocorrelation
signatures (Diebold, 2021). No proper long-memory model (e.g.\
ARFIMA) is benchmarked, so the mechanism is named but not isolated,
and the persistence of fit is conflated with the process behind it.
Second, beyond the loss-choice flag above, QLIKE's finite-sample
edge under noisy proxies matters most where vol forecasts feed
downstream (VaR scaling under Basel III, option-implied surfaces,
dynamic delta hedging), exactly the deployment cases where CSV's
MSE-primary ranking can mislead. Third, the headline NN-vs-HAR
comparison rests on asymmetric ensembling: NN10 averages the top 10
of 100 random seeds, while RF and Bagging receive only within-model
variance reduction (across trees) with no analogous seed or
bootstrap ensembling at the model-class level. A symmetric protocol
(re-fitting RF and Bagging across multiple bootstrap initialisations
and averaging the top-K) could narrow CSV's M\_HAR NN advantage.
Fourth, the 2001--2017 test window mixes a financial crisis, the
European sovereign-debt episode and the pre-COVID calm. Whether the
long-horizon ML gain survives a different regime composition is an
empirical question. The present sample flips its direction, which I
address by isolating regime, cross-section size and information set
as competing explanations (Section~3).

\lead{Replication scope.} All five paper findings are tested on
three DJIA stocks (AAPL, AMZN, JPM), 2016--2024, with all 22 paper
models on both information sets. Seven of CSV's nine M\_ALL
covariates are sourced from public data (HSI, ADS, US3M, EPU, EA,
M1W, \$VOL); per-stock IV (licensed) is proxied by VIX, which
replaces the paper's separate IV and VIX inputs. Being index-level,
VIX underweights stock-specific risk premia. CSV's M\_ALL headline
gains also depend on OptionMetrics IV (a licensed commercial
dataset), meaning the practical ``M\_ALL ML advantage'' is available
only to OptionMetrics subscribers, a deployment constraint CSV do
not discuss. M\_ALL therefore has 11 features (3 HAR lags + 8
covariates) against the paper's 12. M1W and \$VOL are constructed
from OHLCV. Headline tables and figures use the eight models most
relevant to the findings reported here; the 22-model contest at
$h=22$ across both information sets, VaR diagnostics, ALE weights
and MCS retention sit in the appendix (Tables~\ref{tab:var}--\ref{tab:mcs}), with the full
multi-horizon contest in supplementary CSVs. The replication finds
CSV's long-horizon ML advantage depends on the volatility regime of
the test window, while the persistence pattern documented in their
Figure~8 replicates on the present sample.

% =====================================================================
\section{Implementation and headline results}
% =====================================================================

\lead{Data and RV construction.} 1-minute OHLCV data are used for
AAPL, AMZN and JPM, 2016--2024 (2{,}264 NYSE sessions, 390 bars/day).
Following Andersen \& Bollerslev (1998), daily RV is computed as the
sum of squared 5-minute log-returns within day (overnight excluded);
realised quarticity (Bollerslev, Patton \& Quaedvlieg, 2016), signed
semivariance (Patton \& Sheppard, 2015) and aggregated negative
returns (Corsi \& Renò, 2012) support LevHAR, SHAR and HARQ.
Descriptive statistics are in Table~\ref{tab:desc}. Figure~\ref{fig:rv} plots
the unconditional RV path and shows the COVID-19 spike of March
2020 and the August 2024 carry-trade unwind.

\begin{figure}[!htbp]
\centering
\includegraphics[width=\columnwidth]{fig1_rv_overview.png}
\caption{Daily 5-min realised volatility (annualised, \%) for AAPL,
AMZN and JPM, 2016--2024. COVID-19 (Mar 2020) and the Aug 2024
carry-trade unwind are visible. Test window: Mar 2023 -- Dec 2024,
mean $\sigma\approx18\,\%$ vs full-sample 20\,\%.}
\label{fig:rv}
\end{figure}

\lead{Models.} Eight headline models are estimated. The five HAR
variants are fit with OLS: HAR; LogHAR with Jensen correction
$\exp(\hat{y}+0.5\hat{\sigma}^2)$; LevHAR with aggregated
negative-return regressors; SHAR with positive/negative semivariance
decomposition; and HARQ, which conditions the daily-lag coefficient
on realised quarticity via an RVD\,$\times\sqrt{\text{RQ}}$
interaction. Elastic Net (Zou \& Hastie, 2005) hyperparameters are
tuned by grid search on the held-out 10\,\% validation set (k-fold CV
is avoided due to autocorrelation-induced leakage). Random Forest
(Breiman, 2001) uses 500 trees, min-leaf 5, $\sqrt{p}$ random
features per split. The neural network has two hidden layers of 4
and 2 ReLU neurons trained with Adam (Kingma \& Ba, 2014;
lr~0.001); the headline ensemble averages the top 3 of 10 random
seeds ranked by validation MSE (the 22-model appendix contest uses
the paper-exact top 10 of 100; see Deviations). Features are
standardised on the training sample; negative forecasts are floored
at the in-sample RV minimum.

\lead{Forecast targets and evaluation.} Forecasts are produced at
$h\in\{1,5,22\}$ trading days; the target is the $h$-day mean future
RV. With the 70/10/20 split, the test window is March 2023 --
December 2024. Out-of-sample MSE relative to HAR is reported as the
headline metric for like-for-like comparability with CSV Table~2;
QLIKE (Table~\ref{tab:qlike}) is treated as the operationally relevant ranking in
Section~3. Diebold--Mariano (1995) tests run on the squared-error
differentials, using the Harvey, Leybourne \& Newbold (1997)
small-sample correction (which compensates for the variance
inflation induced by overlap in multi-step forecasts) and a
Newey--West long-run variance with $h-1$ Bartlett lags. The full
numerical results sit in Table~\ref{tab:main}; Figure~\ref{fig:relmse}
visualises them. The MCS
uses moving-block bootstrap with block length
$L=\lceil T^{1/3}\rceil$ and 1{,}000 replications. ALE (Apley \&
Zhu, 2020) uses 20 quantile bins; it is preferred over PDP because
RVD/RVW/RVM are highly correlated, breaking PDP's independence
assumption. VaR uses filtered historical simulation (Barone-Adesi,
Giannopoulos \& Vosper, 1999) at $\alpha=5\,\%$.

\lead{Deviations from CSV (transparency).} Four choices differ from
the paper: (i) standard ReLU instead of Leaky-ReLU (Maas, Hannun \&
Ng, 2013; sklearn limitation; paper Footnote~7 confirms negligible
effect); (ii) the insanity filter has a universal floor
(paper-faithful) and adds a universal $2\times$ in-sample-max
ceiling (an extension introduced here); the HARQ-specific
Bollerslev-Patton-Quaedvlieg (2016) filter is omitted; (iii)
HAR-family models are re-estimated daily-rolling in the 22-model
contest (Tables~\ref{tab:var}--\ref{tab:mcs}, paper-faithful) but fixed-window in the
8-model headline contest (Tables~\ref{tab:desc}--\ref{tab:qlike} and headline figures) for compute
reasons; other models use fixed-window throughout, whereas the
paper's Section~2 rolls daily; (iv) NN ensembling is paper-exact
(top 10 of 100) in the 22-model contest, simplified to top-3-of-10
in the 8-model contest; drop-out, learning-rate shrinkage and early
stopping are not implemented; (v) for compute, four grids/specs are
coarsened: GB depth $\in\{3,5\}$, trees $\in\{100,300,500\}$,
lr $\in\{0.01,0.05,0.1\}$ (paper Table~A.6: depth $\{1,2\}$, trees
$\{50,\dots,500\}$, lr $\{0.01,0.1\}$); regularised-linear $\alpha$
grid is 100 log-spaced points on $[10^{-5},10^{2}]$ (paper finer);
RF \texttt{max\_features}\,$=\sqrt{p}$ (paper $J/3$, equivalent at
$J=3$ in M\_HAR, mild divergence at $J=11$ in M\_ALL); ALE uses 20
quantile bins (paper 100, p.~22). Qualitative impact expected small.

\lead{Headline results.} At $h=1$, LogHAR is the strongest
HAR-family member on average (rel-MSE 0.950, DM-significant at
1\,\% for JPM and 5\,\% for AAPL); the neural network matches on
AAPL ($0.934^{**}$), as CSV report for M\_HAR. LogHAR's edge is
largest at $h=5$ (0.856) and persists at $h=22$ (0.887). Random
Forest moves the other way: monthly rel-MSE averages 2.077, driven
by AMZN (3.01) and JPM (2.02). Results vary by stock: on JPM at
$h=5$, RF beats HAR significantly ($0.871^{*}$), so the
cross-section average hides stock-level patterns analysed in
Section~3. Elastic Net offers no improvement under M\_HAR (rel-MSE
0.99--1.07, worst on JPM at 1.22, $h=5$), matching the paper's
claim that ML gains arise mostly under M\_ALL.
Figure~\ref{fig:fc} confirms this on AAPL.

% =====================================================================
\section{Critical comparison and discussion}
% =====================================================================

\lead{What replicates.} Four of CSV's five findings replicate;
Finding~(iv) splits: its persistence mechanism replicates, its MSE
headline does not (addressed below). (a) NN beats HAR on AAPL at
$h=1$ ($0.934^{**}$) and $h=22$ ($0.803^{**}$), supporting
Finding~(i); LogHAR's MSE drops to $\sim0.45$ of HAR's for JPM and
0.69 for AAPL in the upper-middle RV deciles
(Figure~\ref{fig:decile}).
(b)~Under M\_ALL, regularisation and LogHAR cut MSE by 10--30\,\% at
long horizons (Table~\ref{tab:mhar-mall}), replicating Finding~(ii) for regularised
linears and LogHAR but not tree ensembles (see below). (c)~The
persistence pattern replicates: at $h=22$ the ACF of fitted RV
decays much more slowly for RF and NN than for HAR
(Figure~\ref{fig:acf}). For AAPL at lag 60, HAR's ACF is $<0.10$,
RF's $\approx0.22$, NN's $\approx0.15$. AMZN doesn't fit the
pattern: NN decays at HAR's rate, so the persistence is RF-specific
on that stock. To probe the mechanism Diebold (2021) flags, I fit
an ARFIMA$(1,\hat{d},1)$ with $\hat{d}$ from Geweke-Porter-Hudak
(1983): $\hat{d}=0.48$ (AAPL), 0.49 (AMZN), 0.43 (JPM), all in the
stationary long-memory region. Rel-MSE at $h=22$ vs HAR is
$0.72/1.27/0.57$ (mean 0.85). The proper long-memory model beats
HAR; ML's persistent fit (Figure~\ref{fig:acf}) does not. The
plausible reading: ARFIMA exploits the slow-decaying autocorrelation
structure directly, HAR approximates it parsimoniously with three
lag horizons, while RF/NN fit high-frequency nonlinearity on top of
the HAR lags rather than extending the long-memory reach; their
persistent fit (Figure~\ref{fig:acf}) is a consequence of averaging,
not of capturing additional long-tail mass.

\begin{figure}[!htbp]
\centering
\includegraphics[width=\columnwidth]{fig2_acf_persistence.png}
\caption{Reproduction of CSV's Figure~8. ACF of fitted RV at
$h=1$ and $h=22$ for HAR, LogHAR, RF, NN. At $h=22$ the ML models
retain autocorrelation beyond lag~60 while HAR decays to zero. An
ARFIMA$(1,\hat{d},1)$ benchmark on the same series
($\hat{d}\approx0.47$, rel-MSE 0.85 at $h=22$ vs HAR) shows the
long-memory mechanism is real in the RV process, but ML's
persistent fit does not exploit it more efficiently than HAR.}
\label{fig:acf}
\end{figure}

The mechanism is real in the underlying process, but ML's M\_ALL
gain is not a more efficient long-memory approximation than
HAR/LogHAR. (d)~VaR replicates Finding~(v): 21 of 22 models pass
Kupiec (1995) at 5\,\% (Adaptive Lasso fails, $p=0.005$); all 22
pass Christoffersen (1998) (Table~\ref{tab:var}). (e)~Finding~(iii) replicates:
every model except RF places RVD first (RF puts RVW marginally
above RVD: 0.46 vs 0.43); LogHAR is the only model to give the
weekly and monthly lags non-trivial weight (Table~\ref{tab:ale}). Apley \& Zhu
(2020) caveat: the global ranking is a derived stdev-of-ALE
construction, not validated by their framework.

\lead{What does not replicate.} Finding~(iv)'s MSE headline (gains
rising with horizon) is information-set-conditional in the present
sample. Under M\_HAR the headline inverts: RF rel-MSE rises from
1.160 at $h=1$ to 2.077 at $h=22$ on average, against CSV's
RF~$\approx1.02$ at $h=22$ on M\_HAR. Under M\_ALL the paper's
headline is recovered (point~(b); Table~\ref{tab:mhar-mall}). I expected the gain to
weaken but stay directional; instead it inverts. Four explanations
are plausible. (1)~Regime: the paper's window (mid-2014 to end-2017)
covers Aug-2015 China, Brexit and the 2016 election; the present
window (Mar-2023 to Dec-2024) has cross-section mean annualised
$\sigma$ of 18.1\,\% against the full-sample mean of 20.0\,\%, at
the 19th percentile of rolling 21-month windows since 2016
(Figure~\ref{fig:rv}). (2)~Cross-section: with 3 stocks the standard error
scales by $\sqrt{29/3}\approx3.1\times$ under equal-per-stock noise,
so my 2.077 carries roughly $3\times$ the sampling noise of CSV's
1.02. With only 3 stocks the cross-section uncertainty is best
captured by enumerating the three 2-of-3 pairs (AAPL+JPM: 1.61,
AAPL+AMZN: 2.10, AMZN+JPM: 2.52); a $B=1{,}000$ with-replacement
bootstrap returns 95\,\% CI [1.20, 3.01], degenerate to [min, max]
given the 3-point support. Either reading confirms AMZN as the
single driver. CSV's DJIA panel is also index-curated, not a random
large-cap sample: AAPL joined the index in 2015 and AMZN in 2020,
so their 2001--2017 histories enter the CSV panel before they would
have been selected ex ante. If the panel were restricted to ex ante
DJIA-eligible firms, the long-horizon RF gain might attenuate.
(3)~Information set: M\_ALL is where CSV's main ML gains arise
(Table~\ref{tab:mhar-mall}). (4)~Rolling refit narrows AMZN by $\sim13\,\%$ but the
qualitative picture is unchanged (Figure~\ref{fig:rolling}). Secondary: LogHAR's
dominance reflects the calm window favouring outlier
down-weighting, and tree ensembles deteriorate on M\_ALL (Bagging,
GBoost; Table~\ref{tab:mhar-mall}), plausibly because 3 stocks concentrate
overfitting that 29 would smooth out.

\lead{Loss-function robustness.} Patton (2011) shows that MSE and
QLIKE both preserve the expected ranking of $\sigma^2$-targeting
forecasts under noisy proxies, but QLIKE produces loss differences
with lower finite-sample variance and higher pairwise-test power
(QLIKE loss $\sigma^2/\hat{\sigma}^2 - \log(\sigma^2/\hat{\sigma}^2)
- 1$). Table~\ref{tab:qlike} re-evaluates under QLIKE. LogHAR remains best at
$h=22$ (rel-QLIKE $0.831^{*}$); RF's monthly deterioration shrinks
from 2.077 (MSE) to 1.279 (QLIKE); on AAPL at $h=22$ RF beats HAR
under QLIKE (0.870 vs 1.201 under MSE). HARQ shows the largest
HAR-family MSE-to-QLIKE shift ($1.210\to0.957$), consistent with BPQ
(2016): their RQ-conditioned daily-lag coefficient down-weights
observations where measurement noise drives the daily lag away from
its long-run value, and QLIKE assigns lower weight to those same
observations. The AMZN deterioration survives QLIKE (1.844). Read
with Table~\ref{tab:per-stock-qlike}, the picture sharpens: under proxy-aware losses the
long-horizon ML disadvantage is concentrated on AMZN, so the
cross-section generalisation question is whether the AMZN pattern
persists on a larger panel. For risk-management uses (VaR scaling,
option-implied surfaces), QLIKE is the operationally relevant
ranking.

\lead{Insight.} Under a proxy-aware loss the long-horizon ML
failure reduces to AMZN. The ARFIMA$(1,\hat{d},1)$ benchmark
($\hat{d}\approx0.47$, rel-MSE 0.85 at $h=22$) shows the
long-memory mechanism is genuine; ML doesn't exploit it more
efficiently than HAR. CSV report MSE primary and QLIKE only as
robustness; the present sample shows that ordering can flip the
headline. For risk-management uses (VaR scaling, option-implied
surfaces) QLIKE is the operationally relevant loss, and the AMZN
inversion is the only signal that survives it. The MCS reinforces
this on M\_HAR: at the 90\,\% level the HAR family is retained in 7
or more of 9 cells while Adaptive Lasso is eliminated in all 9 at
the 75\,\% level --- uniformly across AAPL, AMZN and JPM, so the
parsimony advantage reflects general small-cross-section
instability of two-stage shrinkage rather than AMZN-specific
outlier protection. The AMZN-driver evidence sits in Table~\ref{tab:per-stock-qlike}'s
per-stock QLIKE view: under MSE the failure pattern looks like
regime drift, under QLIKE it collapses to AMZN. Basel~III VaR
scaling, variance-swap valuation and dynamic delta hedging depend on
persistent-shock forecasts stressed in CSV's 2014--2017 window but
not the present one --- the deployment cases where ML's edge would
matter most are the ones this sample cannot probe. On the modelling
side, features dominate depth: deeper NNs do not beat shallow ones,
and single-fit NNs are high-variance ($\mathrm{NN}_3^{1}$ at $h=22$
$=1.48$ vs $\mathrm{NN}_3^{10}=1.04$; Table~\ref{tab:mhar-mall}), supporting the
paper's ensembling emphasis. Branco, Rubesam \& Zevallos (2024) find
that nonlinear ML does not statistically beat HAR-class linear
models on ten global indices, consistent with the M\_HAR result
here. A simpler benchmark goes further. GARCH(1,1) (Bollerslev 1986)
fit on daily returns forecasts return variance, equivalent to
integrated variance under no-jump (Andersen \& Bollerslev, 1998),
and beats HAR cross-section at $h=22$ (rel-MSE 0.66; AAPL 0.54,
AMZN 0.98, JPM 0.45), outperforming ARFIMA's 0.85. Hansen \& Lunde
(2005) framed GARCH(1,1) as the canonical pre-ML benchmark. So the
HAR baseline is itself beatable by simple alternatives, and ML's
M\_HAR failure is sharpened: ML's persistent fit doesn't even reach
the level of a 1-equation GARCH. CSV's
$22\times29\times3\times2=3{,}828$ pairwise DM tests imply an
expected $\sim$191 false rejections at $\alpha=0.05$ without
correction; MCS controls error within each (stock, horizon) cell but
not across cells. I would update toward CSV's structural reading if
their long-horizon RF gain held on a calm post-2017 sample with
three to five stocks matching mine, or if rolling re-estimation with
stock-by-stock feature engineering removed the AMZN-specific
deterioration.

% =====================================================================
%   REFERENCES  (do not count toward 3-page body limit)
% =====================================================================
\onecolumn
\section*{References}
\small
\setlength{\parskip}{1pt}

\begingroup
\setlength{\parindent}{-1em}\leftskip1em\relax

Andersen, T. G., \& Bollerslev, T. (1998). Answering the skeptics:
Yes, standard volatility models do provide accurate forecasts.
\textit{International Economic Review}, 39(4), 885--905.\par
Apley, D. W., \& Zhu, J. (2020). Visualizing the effects of
predictor variables in black box supervised learning models.
\textit{Journal of the Royal Statistical Society: Series B}, 82(4),
1059--1086.\par
Baker, S. R., Bloom, N., \& Davis, S. J. (2016). Measuring economic
policy uncertainty. \textit{The Quarterly Journal of Economics},
131(4), 1593--1636.\par
Barone-Adesi, G., Giannopoulos, K., \& Vosper, L. (1999). VaR
without correlations for portfolios of derivative securities.
\textit{Journal of Futures Markets}, 19(5), 583--602.\par
Bollerslev, T. (1986). Generalized autoregressive conditional
heteroskedasticity. \textit{Journal of Econometrics}, 31(3),
307--327.\par
Bollerslev, T., Patton, A. J., \& Quaedvlieg, R. (2016). Exploiting
the errors: A simple approach for improved volatility forecasting.
\textit{Journal of Econometrics}, 192(1), 1--18.\par
Branco, R. R., Rubesam, A., \& Zevallos, M. (2024). Forecasting
realized volatility: Does anything beat linear models?
\textit{Journal of Empirical Finance}, 78, 101524.\par
Breiman, L. (1996). Bagging predictors. \textit{Machine Learning},
24(2), 123--140.\par
Breiman, L. (2001). Random forests. \textit{Machine Learning},
45(1), 5--32.\par
Christensen, K., Siggaard, M., \& Veliyev, B. (2023). A machine
learning approach to volatility forecasting. \textit{Journal of
Financial Econometrics}, 21(5), 1680--1727.\par
Christoffersen, P. F. (1998). Evaluating interval forecasts.
\textit{International Economic Review}, 39(4), 841--862.\par
Corsi, F. (2009). A simple approximate long-memory model of
realized volatility. \textit{Journal of Financial Econometrics},
7(2), 174--196.\par
Corsi, F., \& Renò, R. (2012). Discrete-time volatility forecasting
with persistent leverage effect and the link with continuous-time
volatility modeling. \textit{Journal of Business \& Economic
Statistics}, 30(3), 368--380.\par
Diebold, F. X. (2021). Machine learning for realized volatility
forecasting. \textit{No Hesitations} [blog], 1 February.
\url{https://fxdiebold.blogspot.com/2021/02/machine-learning-for-realized.html}\par
Diebold, F. X., \& Mariano, R. S. (1995). Comparing predictive
accuracy. \textit{Journal of Business \& Economic Statistics},
13(3), 253--263.\par
Friedman, J. H. (2001). Greedy function approximation: A gradient
boosting machine. \textit{Annals of Statistics}, 29(5),
1189--1232.\par
Geweke, J., \& Porter-Hudak, S. (1983). The estimation and
application of long memory time series models. \textit{Journal of
Time Series Analysis}, 4(4), 221--238.\par
Hansen, P. R., \& Lunde, A. (2005). A forecast comparison of
volatility models: Does anything beat a GARCH(1,1)?
\textit{Journal of Applied Econometrics}, 20(7), 873--889.\par
Hansen, P. R., Lunde, A., \& Nason, J. M. (2011). The Model
Confidence Set. \textit{Econometrica}, 79(2), 453--497.\par
Harvey, D., Leybourne, S., \& Newbold, P. (1997). Testing the
equality of prediction mean squared errors. \textit{International
Journal of Forecasting}, 13(2), 281--291.\par
Kingma, D. P., \& Ba, J. (2014). Adam: A method for stochastic
optimization. \textit{arXiv:1412.6980}.\par
Kupiec, P. H. (1995). Techniques for verifying the accuracy of risk
measurement models. \textit{Journal of Derivatives}, 3(2),
73--84.\par
Maas, A. L., Hannun, A. Y., \& Ng, A. Y. (2013). Rectifier
nonlinearities improve neural network acoustic models.
\textit{Proceedings of the 30th International Conference on Machine
Learning}, 30(1), 3.\par
Patton, A. J. (2011). Volatility forecast comparison using
imperfect volatility proxies. \textit{Journal of Econometrics},
160(1), 246--256.\par
Patton, A. J., \& Sheppard, K. (2015). Good volatility, bad
volatility: Signed jumps and the persistence of volatility.
\textit{Review of Economics and Statistics}, 97(3), 683--697.\par
Tibshirani, R. (1996). Regression shrinkage and selection via the
lasso. \textit{Journal of the Royal Statistical Society: Series B},
58(1), 267--288.\par
Zou, H., \& Hastie, T. (2005). Regularization and variable
selection via the elastic net. \textit{Journal of the Royal
Statistical Society: Series B}, 67(2), 301--320.\par
\endgroup

% =====================================================================
%   APPENDIX — tables and figures (do not count toward 3-page body)
% =====================================================================
\clearpage
\appendix
% Tight float spacing for the appendix so tables/figures pack closely.
\setlength{\floatsep}{4pt plus 1pt minus 1pt}
\setlength{\textfloatsep}{6pt plus 1pt minus 1pt}
\setlength{\intextsep}{4pt plus 1pt minus 1pt}
\setlength{\abovecaptionskip}{2pt}
\setlength{\belowcaptionskip}{2pt}

\section*{Appendix A: Tables}

% ------------------ Table 1 ------------------
\begin{table}[H]
\centering\small
\caption{Descriptive statistics of daily realised variance
(5-minute subsampled). RV is reported scaled by $10^{4}$; the
annualised $\sigma$ column is the in-sample mean of
$\sigma_t=\sqrt{252\cdot\mathrm{RV}_t}$.}
\label{tab:desc}
\begin{tabular}{lllrrrrrrr}
\toprule
Ticker & Sector & N days & Start & End & Mean RV & Median RV & Ann.\ $\sigma$ (\%) & Skew & Kurt \\
\midrule
AAPL & Technology & 2{,}264 & 2016-01-04 & 2024-12-31 & 1.89 & 1.14 & 19.3 & 8.65 & 119.24 \\
AMZN & Consumer   & 2{,}264 & 2016-01-04 & 2024-12-31 & 2.53 & 1.54 & 22.4 & 6.53 & 87.67 \\
JPM  & Financials & 2{,}264 & 2016-01-04 & 2024-12-31 & 1.76 & 1.04 & 18.4 & 12.23 & 190.65 \\
\bottomrule
\end{tabular}
\end{table}

% ------------------ Table 2 ------------------
\begin{table}[H]
\centering\small
\caption{Out-of-sample MSE relative to HAR. Lower is better;
1.000 = HAR baseline. Stars indicate Diebold--Mariano rejection of
equal predictive accuracy in favour of the row model, with
Harvey--Leybourne--Newbold small-sample correction (*:~10\,\%,
**:~5\,\%, ***:~1\,\%). EN = Elastic Net; RF = Random Forest;
NN = Neural Network (4, 2).}
\label{tab:main}
\begin{tabular}{llrrrr}
\toprule
Horizon & Model & AAPL & AMZN & JPM & Mean \\
\midrule
\multirow{8}{*}{$h=1$}
 & HAR     & 1.000      & 1.000  & 1.000      & 1.000 \\
 & LogHAR  & 0.934**    & 1.003  & 0.913***   & 0.950 \\
 & LevHAR  & 1.013      & 1.053  & 1.037      & 1.034 \\
 & SHAR    & 0.976      & 0.977  & 1.195      & 1.049 \\
 & HARQ    & 0.982      & 0.989  & 1.133      & 1.035 \\
 & EN      & 1.003      & 1.000  & 0.980*     & 0.995 \\
 & RF      & 1.179      & 1.295  & 1.006      & 1.160 \\
 & NN      & 0.934**    & 1.056  & 0.948**    & 0.979 \\
\midrule
\multirow{8}{*}{$h=5$}
 & HAR     & 1.000      & 1.000  & 1.000      & 1.000 \\
 & LogHAR  & 0.816*     & 1.066  & 0.685***   & 0.856 \\
 & LevHAR  & 0.983      & 0.978  & 1.065      & 1.009 \\
 & SHAR    & 1.041      & 0.997  & 1.071      & 1.036 \\
 & HARQ    & 0.997      & 1.062  & 1.480      & 1.180 \\
 & EN      & 0.987      & 0.999  & 1.221      & 1.069 \\
 & RF      & 1.071      & 1.659  & 0.871*     & 1.200 \\
 & NN      & 0.906*     & 1.284  & 0.955      & 1.049 \\
\midrule
\multirow{8}{*}{$h=22$}
 & HAR     & 1.000      & 1.000  & 1.000      & 1.000 \\
 & LogHAR  & 0.686***   & 1.353  & 0.623***   & 0.887 \\
 & LevHAR  & 0.849***   & 0.983  & 1.161      & 0.998 \\
 & SHAR    & 1.078      & 0.994  & 1.007      & 1.027 \\
 & HARQ    & 1.006      & 1.083  & 1.540      & 1.210 \\
 & EN      & 0.998      & 1.000  & 1.203      & 1.067 \\
 & RF      & 1.201      & 3.007  & 2.023      & 2.077 \\
 & NN      & 0.803**    & 1.501  & 1.128      & 1.144 \\
\bottomrule
\end{tabular}
\end{table}

% ------------------ Table 3 ------------------
\begin{table}[H]
\centering\small
\caption{QLIKE robustness check (Patton, 2011). Cross-section mean
of out-of-sample QLIKE relative to HAR; lower is better. A `*'
indicates DM rejection on the QLIKE differential in at least 2 of 3
stocks at the 10\,\% level. Numbers cited in Section~3
(Loss-function robustness, Insight) are reproduced here for direct
verification.}
\label{tab:qlike}
\begin{tabular}{lrrrrrrr}
\toprule
Horizon & LogHAR & LevHAR & SHAR & HARQ & EN & RF & NN \\
\midrule
$h=1$   & 0.961   & 1.060 & 1.154 & 1.057 & 1.031 & 1.069 & 0.977* \\
$h=5$   & 0.945*  & 0.994 & 1.018 & 1.059 & 1.185 & 1.093 & 1.121  \\
$h=22$  & 0.831*  & 0.879 & 1.003 & 0.957 & 1.084 & 1.279 & 1.040  \\
\bottomrule
\end{tabular}
\end{table}

% ------------------ Table 3b ------------------
% Numbering trick: force this table to show as "3b" without advancing the
% counter, so subsequent tables remain 4, 5, 6, 7 as written in the body.
\renewcommand{\thetable}{3b}
\addtocounter{table}{-1}
\begin{table}[H]
\centering\small
\caption{Per-stock comparison of MSE vs QLIKE at $h=22$ (rel-HAR).
For risk-management use cases (VaR scaling, variance swaps) QLIKE is
the operationally relevant loss. Under QLIKE the long-horizon ML
deterioration concentrates on AMZN — RF beats HAR on both AAPL and
JPM (0.870, ---) but blows up on AMZN (1.844).}
\label{tab:per-stock-qlike}
\begin{tabular}{l|rr|rr|rr}
\toprule
       & \multicolumn{2}{c|}{AAPL} & \multicolumn{2}{c|}{AMZN}
       & \multicolumn{2}{c}{JPM} \\
Model  & MSE & QLIKE & MSE & QLIKE & MSE & QLIKE \\
\midrule
LogHAR        & 0.686 & 0.724 & 1.353 & 1.141 & 0.623 & 0.627 \\
LevHAR        & 0.849 & 0.845 & 0.983 & 0.940 & 1.161 & 0.852 \\
SHAR          & 1.078 & 1.009 & 0.994 & 0.995 & 1.007 & 1.006 \\
HARQ          & 1.006 & 0.963 & 1.083 & 1.018 & 1.540 & 0.890 \\
EN            & 0.998 & 1.003 & 1.000 & 1.003 & 1.203 & 1.246 \\
RF            & 1.201 & 0.870 & 3.007 & 1.844 & 2.023 & 1.123 \\
NN(4,2)       & 0.803 & 0.783 & 1.501 & 1.225 & 1.128 & 1.113 \\
\bottomrule
\end{tabular}
\end{table}
% Restore default table numbering (\arabic{table}) for tables 4 onward.
\renewcommand{\thetable}{\arabic{table}}

% ------------------ Table 4 ------------------
\begin{table}[H]
\centering\small
\caption{Value-at-Risk diagnostics for the 22-model contest on
M\_HAR (one-day-ahead, $\alpha=5\,\%$). Hit rate, Kupiec
unconditional-coverage $p$-value, and Christoffersen
conditional-independence $p$-value, cross-section mean across the
three stocks. Twenty-one of 22 models pass Kupiec at the 5\,\%
level; only Adaptive Lasso under-covers significantly
($p=0.005$). All 22 models pass Christoffersen.}
\label{tab:var}
\begin{tabular}{lrrr}
\toprule
Model & Hit rate & $p$ (Kupiec) & $p$ (Christoffersen) \\
\midrule
HAR         & 0.039 & 0.283 & 0.452 \\
HAR-X       & 0.039 & 0.283 & 0.452 \\
LogHAR      & 0.045 & 0.605 & 0.661 \\
LevHAR      & 0.045 & 0.516 & 0.172 \\
SHAR        & 0.041 & 0.396 & 0.278 \\
HARQ        & 0.045 & 0.408 & 0.383 \\
Ridge       & 0.037 & 0.229 & 0.436 \\
Lasso       & 0.035 & 0.184 & 0.470 \\
ElasticNet  & 0.036 & 0.222 & 0.448 \\
\textbf{AdaLasso}   & \textbf{0.021} & \textbf{0.005} & 0.501 \\
PostLasso   & 0.039 & 0.332 & 0.478 \\
Bagging     & 0.041 & 0.427 & 0.496 \\
RandomForest& 0.044 & 0.442 & 0.418 \\
GBoost      & 0.035 & 0.223 & 0.510 \\
NN$_1^{1}$  & 0.040 & 0.327 & 0.431 \\
NN$_1^{10}$ & 0.041 & 0.427 & 0.486 \\
NN$_2^{1}$  & 0.037 & 0.229 & 0.436 \\
NN$_2^{10}$ & 0.042 & 0.453 & 0.476 \\
NN$_3^{1}$  & 0.027 & 0.077 & 0.542 \\
NN$_3^{10}$ & 0.040 & 0.327 & 0.431 \\
NN$_4^{1}$  & 0.042 & 0.429 & 0.660 \\
NN$_4^{10}$ & 0.041 & 0.376 & 0.458 \\
\bottomrule
\end{tabular}
\end{table}

% ------------------ Table 5 ------------------
\begin{table}[H]
\centering\small
\caption{ALE-based feature importance at $h=1$ (cross-section mean
share over RVD/RVW/RVM, sums to 1.00 by row).  Every model except
RF places RVD first; LogHAR is the only model to put non-trivial
weight on the weekly and monthly lags.}
\label{tab:ale}
\begin{tabular}{lrrr}
\toprule
Model        & RVD  & RVW  & RVM \\
\midrule
HAR          & 0.55 & 0.43 & 0.02 \\
LogHAR       & 0.54 & 0.32 & 0.14 \\
ElasticNet   & 0.54 & 0.46 & 0.00 \\
RandomForest & 0.43 & 0.46 & 0.11 \\
NN(4,2)      & 0.49 & 0.40 & 0.11 \\
\bottomrule
\end{tabular}
\end{table}

% ------------------ Table 6 ------------------
\begin{table}[H]
\centering\scriptsize
\caption{M\_HAR vs M\_ALL: cross-section mean rel-MSE across AAPL,
AMZN, JPM for all 22 models on both information sets at three
horizons. ``RR/LA/EN/A-LA/P-LA'' = Ridge / Lasso / ElasticNet /
Adaptive-Lasso / Post-Lasso; ``BG/RF/GB'' = Bagging / Random Forest /
Gradient Boosting; NN$_k^{n}$ denotes a $k$-hidden-layer net with
top-$n$-of-100 ensembling.  Under M\_ALL most regularised linears
and the NN10 ensembles improve, while LevHAR/SHAR/HARQ stay near
parity and tree ensembles deteriorate at $h=22$.}
\label{tab:mhar-mall}
\begin{tabular}{ll|rr|rr|rr}
\toprule
 & & \multicolumn{2}{c|}{$h=1$} & \multicolumn{2}{c|}{$h=5$} & \multicolumn{2}{c}{$h=22$} \\
Group & Model & M\_HAR & M\_ALL & M\_HAR & M\_ALL & M\_HAR & M\_ALL \\
\midrule
HAR family   & HAR          & 1.000 & 1.000 & 1.000 & 1.000 & 1.000 & 1.000 \\
             & HAR-X        & 1.000 & 1.000 & 1.000 & 1.000 & 1.000 & 1.000 \\
             & LogHAR       & 0.949 & 0.853 & 0.852 & 0.679 & 0.889 & 0.682 \\
             & LevHAR       & 1.030 & 1.008 & 1.010 & 0.978 & 1.008 & 1.021 \\
             & SHAR         & 1.045 & 1.047 & 1.034 & 1.028 & 1.026 & 1.022 \\
             & HARQ         & 1.034 & 1.014 & 1.174 & 1.068 & 1.224 & 1.074 \\
\midrule
Reg.\ linear & RR           & 0.992 & 0.992 & 0.995 & 1.004 & 1.077 & 1.078 \\
             & LA           & 0.998 & 0.990 & 1.091 & 0.872 & 1.167 & 0.847 \\
             & EN           & 0.997 & 1.016 & 1.089 & 0.908 & 1.158 & 0.899 \\
             & A-LA         & 1.669 & 1.724 & 1.722 & 1.512 & 1.755 & 1.327 \\
             & P-LA         & 0.996 & 1.033 & 1.000 & 0.907 & 1.073 & 0.817 \\
\midrule
Tree         & BG           & 1.291 & 1.076 & 1.396 & 1.023 & 3.024 & 3.988 \\
             & RF           & 1.162 & 0.979 & 1.218 & 0.914 & 2.233 & 2.330 \\
             & GB           & 1.245 & 1.190 & 1.164 & 1.056 & 1.876 & 4.533 \\
\midrule
Neural net   & NN$_1^{1}$   & 1.009 & 0.996 & 1.001 & 1.085 & 0.996 & 1.526 \\
             & NN$_1^{10}$  & 1.031 & 0.914 & 0.991 & 0.828 & 1.162 & 1.003 \\
             & NN$_2^{1}$   & 0.982 & 1.062 & 1.183 & 1.043 & 1.414 & 1.713 \\
             & NN$_2^{10}$  & 1.007 & 0.911 & 1.252 & 0.773 & 1.358 & 0.999 \\
             & NN$_3^{1}$   & 1.402 & 1.173 & 1.300 & 2.452 & 1.203 & 1.477 \\
             & NN$_3^{10}$  & 1.007 & 0.895 & 1.025 & 0.818 & 1.485 & 1.041 \\
             & NN$_4^{1}$   & 1.037 & 1.283 & 1.187 & 1.526 & 1.592 & 1.894 \\
             & NN$_4^{10}$  & 1.020 & 0.933 & 1.070 & 0.891 & 1.511 & 1.294 \\
\bottomrule
\end{tabular}
\end{table}

% ------------------ Table 7 ------------------
\begin{table}[H]
\centering\small
\caption{Model Confidence Set retention on M\_HAR. Counts are the
number of (stock, horizon) cells (out of 9: 3 stocks $\times$ 3
horizons) in which each model survives elimination at the stated
confidence level. The HAR family is retained almost everywhere;
Adaptive Lasso is eliminated in every 75\,\%-level cell.}
\label{tab:mcs}
\begin{tabular}{lrr|lrr}
\toprule
Model & MCS$_{90}$ & MCS$_{75}$ & Model & MCS$_{90}$ & MCS$_{75}$ \\
\midrule
HAR          & 9 & 8 & PostLasso   & 9 & 8 \\
HAR-X        & 9 & 8 & Bagging     & 6 & 3 \\
LogHAR       & 9 & 8 & RandomForest& 7 & 5 \\
LevHAR       & 8 & 8 & GBoost      & 7 & 4 \\
SHAR         & 8 & 7 & NN$_1^{1}$  & 9 & 8 \\
HARQ         & 7 & 5 & NN$_1^{10}$ & 7 & 6 \\
Ridge        & 9 & 8 & NN$_2^{1}$  & 7 & 6 \\
Lasso        & 8 & 8 & NN$_2^{10}$ & 6 & 6 \\
ElasticNet   & 8 & 8 & NN$_3^{1}$  & 6 & 6 \\
\textbf{AdaLasso}    & \textbf{0} & \textbf{0} & NN$_3^{10}$ & 7 & 5 \\
             &   &   & NN$_4^{1}$  & 7 & 7 \\
             &   &   & NN$_4^{10}$ & 7 & 6 \\
\bottomrule
\end{tabular}
\end{table}

% =====================================================================
%   Appendix B — figures (2x2 grid, single page)
% =====================================================================
\clearpage
\section*{Appendix B: Figures}

% --- Row 1: forecasts + rolling robustness ---
\noindent
\begin{minipage}[t]{0.48\linewidth}
\centering
\includegraphics[width=\linewidth]{fig4_forecasts.png}
\captionof{figure}{AAPL one-day-ahead forecasts vs realised
volatility. HAR, LogHAR and NN are indistinguishable in the centre
of the distribution; differences emerge at volatility spikes
(August 2024).}
\label{fig:fc}
\end{minipage}\hfill
\begin{minipage}[t]{0.48\linewidth}
\centering
\includegraphics[width=\linewidth]{fig6_rolling_robust.png}
\captionof{figure}{Rolling-refit robustness. HAR and RF
re-estimated every 22 days narrow AMZN by $\sim13\,\%$
($3.01\to2.63$), JPM by $\sim4\,\%$, AAPL marginally. Rolling RF
uses 300 trees vs 500 static, a lower bound on the pure refit
effect.}
\label{fig:rolling}
\end{minipage}

\vspace{8pt}

% --- Fig 5 (relmse) and Fig 6 (decile) stacked full-width, same page as Row 1 ---
\begin{figure}[H]
\centering
\includegraphics[width=0.9\textwidth]{fig3_relative_mse.png}
\caption{Out-of-sample MSE relative to HAR for each
model--stock--horizon. Dashed line at 1.0 is the HAR benchmark;
values below 1.0 indicate improvement. LogHAR is the strongest HAR
variant at every horizon; Random Forest deteriorates sharply at
$h=22$, especially on AMZN. $y$-axis range varies across stocks to
accommodate the AMZN $h=22$ RF outlier.}
\label{fig:relmse}
\end{figure}

\begin{figure}[H]
\centering
\includegraphics[width=0.9\textwidth]{fig5_decile_mse.png}
\caption{Reproduction of CSV's Figure~5. Forecast MSE relative to
HAR by decile of realised RV in the test set, $h=1$. LogHAR is
broadly better than HAR for AAPL and JPM, with parity in the middle
deciles; RF and NN are worse in the middle deciles for AMZN but
converge to HAR in the top decile, where LogHAR's
outlier-down-weighting dominates.}
\label{fig:decile}
\end{figure}

\end{document}